\scanpage{127}
\mspage{25}
\deleted{obtenue par extension du corps de base}
\deleted{
obtient ainsi une ``fibration'' naturelle de l'ouvert dense $U$ des points réguliers de $G$ au dessus du \deleted{\unsure{variété}} préschéma $\deleted{\unsure{U}}\ \unsure{\mathcal{C}|\mathcal{T}}$, les fibres \deleted{géométriques} étant \deleted{\illegible{1}} \inserted{des} ouverts \inserted{denses} de sous-groupes de Cartan des fibres de $G_{\deleted{\unsure{U}}\unsure{\mathcal{C}|\mathcal{T}}}$ (savoir les ouverts des points réguliers dans $G$). \deleted{modulo extension du corps de base.} On trouve par exemple\inserted{: le résultat suivant (qui sera considérablement précisé dans l'exposé suivant):}

\uline{Corollaire 3.3.} Soient $G$ un groupe algébrique lisse connexe \deleted{affine} sur le corps $k$, $\unsure{\mathcal{C}|\mathcal{T}}$ le schéma des tores maximaux de $G$ ($\simeq$ le schéma des sous-groupes de Cartan de $G$). Alors le corps des \deleted{fonctions} fonctions $k(G)$ \inserted{de $G$} est isomorphe au corps de fonctions d'un groupe algébrique lisse connexe affine nilpotent $C$ sur le corps des fonctions $k(\unsure{\mathcal{C}|\mathcal{T}})$ de $\unsure{\mathcal{C}|\mathcal{T}}$, savoir $C=\deleted{\unsure{a}}$ ``le sous-groupe de Cartan générique de $G$''. Si \deleted{$k$ est de rang unipotent nul, i.e.} \inserted{$G$ est affine de rang unipotent nul, i.e. si} les sous-groupes de Cartan \inserted{de $G_{\bar{k}}$} sont des tores, alors $k(G)$ est une extension unirationnelle de $k(\unsure{\mathcal{C}|\mathcal{T}})$.

Bien entendu, par sous-groupe de Cartan générique de $G$, on entend (par abus de langage) le sous-groupe de Cartan de $G_{k(\unsure{\mathcal{C}|\mathcal{T}})}$ fibre générique de $X$ au dessus de $\unsure{\mathcal{C}|\mathcal{T}}$. Il n'y a plus qu'à prouver la dernière assertion de 3.3., qui est contenue dans le \inserted{résultat bien connu suivant (dû à Chevalley):}

\deleted{\uline{Lemme 3.4.} Soient $k$ un corps, $T$ un tore sur $k$, $k(T)$ le corps des fonctions rationnelles sur $T$, alors $k(T)$ est une extension unirationnelle de $k$, i.e. est contenue dans une extension transcendante pure de $k$.

Soit en effet $k'$ une extension finie \deleted{\illegible{1}} séparable de $k$ \deleted{\unsure{a}} qui splitte $T$ (Exp X 1.4.), alors \deleted{\unsure{le}} $T\otimes_k k'$ est une variété rationnelle i.e. \deleted{\unsure{isomorphe bir}} admet un ouvert dense isomorphe à un ouvert dense de l'espace euclidien $E_k^n$, donc
\[
T'=\prod_{\operatorname{Spec}(k')/\operatorname{Spec}(k)}T_{k'}/\operatorname{Spec}(k')
\]
est une variété rationnelle (car admet un ouvert dense isomorphe à un ouvert dense de $\prod_{\operatorname{Spec}(k')/\operatorname{Spec}(k)}E_{k'}^n$, qui est isomorphe à l'espace \deleted{\unsure{non}} eucli-
}
}
\margin{Conserver le texte ci-contre, barré par erreur.}
\margin{126}
